\chapter{Reproducibility}
\label{Appendix:reproducibility}

\section{Repository}
\label{sec:appendix-repository}

The complete source code is at \repositorio\ (V.10). Its \texttt{README.md}
reproduces every experiment in this report from a clean machine.

\section{Environment}
\label{sec:appendix-environment}

\todo{Python version, how the virtual environment is created, and the three
pinned requirement files. State that every dependency is pinned with
\texttt{==} and that the exact versions of a given run are recorded in its
\texttt{run.json}.}

\section{Obtaining the data}
\label{sec:appendix-data}

\todo{For each corpus: where it is downloaded from, its size, its checksum, its
licence, and the exact directory the files must end up in. Repeat here what the
README says, so that the report stands on its own.}

\section{One command per experiment}
\label{sec:appendix-commands}

\todo{Reproduce the experiment table of the README: experiment, configuration
file, outputs, and the figure or table of this report it feeds.}

\section{Hardware}
\label{sec:appendix-hardware}

\todo{The laptop the CPU experiments ran on and the department server used for
the trained models and MERT: CPU, RAM, GPU, operating system, and the
wall-clock time of the expensive stages. State plainly which stages need a
GPU (or a slow CPU run) and which do not: the siamese and triplet networks,
the CNN classifier and MERT need a GPU or a slow CPU fallback; the classical
techniques (MFCC, acoustic descriptors, the Gaussian timbre model, the
classical classifiers) and the fuzzy block run on CPU only, in minutes.}
